% Generated by scripts/generate_glossary_latex.jl.
% Edit appendix/GlossaryContent.jl, not this file.
\documentclass[10pt]{article}
\usepackage[a4paper,margin=2.1cm]{geometry}
\usepackage[T1]{fontenc}
\usepackage{lmodern}
\usepackage{microtype}
\usepackage{amsmath,amssymb}
\usepackage{booktabs,longtable,array}
\usepackage{xcolor}
\usepackage{hyperref}
\usepackage{enumitem}
\usepackage{titlesec}
\usepackage{parskip}

\definecolor{mathaccent}{HTML}{3F5F73}
\definecolor{aiaccent}{HTML}{6B5D45}
\definecolor{olive}{HTML}{6A7047}
\definecolor{softgray}{HTML}{F4F4F1}
\definecolor{muted}{HTML}{66665F}

\hypersetup{colorlinks=true,linkcolor=olive,urlcolor=mathaccent,citecolor=olive}
\setlist[itemize]{leftmargin=1.5em,itemsep=0.2em}
\titleformat{\section}{\Large\bfseries}{}{0pt}{}
\titleformat{\subsection}{\large\bfseries\color{olive}}{}{0pt}{}

\newenvironment{glossaryentry}[1]{%
  \par\medskip\noindent
  {\large\bfseries #1}\par\smallskip
}{\par\medskip}

\newenvironment{interpretationnote}[1]{%
  \par\smallskip\noindent
  \begin{minipage}{\linewidth}
  \setlength{\fboxsep}{8pt}%
  \colorbox{softgray}{\begin{minipage}{\dimexpr\linewidth-2\fboxsep}
  \textbf{#1}\par\smallskip
}{%
  \end{minipage}}
  \end{minipage}\par\smallskip
}

\newcommand{\keydistinction}[1]{%
\begin{center}
\fbox{\parbox{0.88\linewidth}{\centering\bfseries #1}}
\end{center}
}

\title{\textbf{Low-Rank Structure Is Geometry}\\
\large Mathematical and AI/Interpretability Glossary}
\author{Companion appendix}
\date{}

\begin{document}
\maketitle

\begin{abstract}
This optional companion groups mathematical and AI/interpretability vocabulary
by topic. Each entry begins with a compact explanation. Only terms with a
substantial pedagogical caution receive a separate \emph{Why it matters} or
\emph{Interpretation boundary} note. Citation keys expand in the final index.
\end{abstract}

\tableofcontents
\newpage

\section{How to use this glossary}
Read the main notebooks first and use this appendix as a lookup reference.
Definitions remain local to the notebooks; this companion provides additional
precision and a bridge between mathematical and interpretability language.

\keydistinction{Coordinates are not automatically intrinsic objects. Learned factors are not automatically validated concepts.}

\section{Mathematical glossary}

\subsection{I. Tensor basics}

\begin{glossaryentry}{Tensor}
A multidimensional array $\mathcal X\in\mathbb R^{n_1\times\cdots\times n_d}$. A matrix is the order-2 case. The notebooks preserve meaningful modes instead of always flattening them.
\par\smallskip{\small\color{muted}[KB09]}
\end{glossaryentry}


\begin{glossaryentry}{Mode}
One axis of a tensor. A mode may represent samples, spatial positions, features, time, or another scientifically meaningful variable class.
\par\smallskip{\small\color{muted}[KB09]}
\end{glossaryentry}


\begin{glossaryentry}{Matricization / unfolding}
A rearrangement of a tensor into a matrix by choosing one or more modes as matrix axes. Flattening can enable matrix methods while merging distinct mode semantics.
\par\smallskip{\small\color{muted}[KB09]}
\end{glossaryentry}


\begin{glossaryentry}{Outer product}
For vectors $a,b,c$, the tensor $a\otimes b\otimes c$ has entries $(a\otimes b\otimes c)_{ijk}=a_i b_j c_k$. Every CP rank-one term is an outer product of one vector from each mode.
\par\smallskip{\small\color{muted}[KB09]}
\end{glossaryentry}


\begin{glossaryentry}{Rank-one tensor}
A tensor that can be written as one outer product $a^{(1)}\otimes\cdots\otimes a^{(d)}$. It is the basic component of CP decomposition and the Segre model.
\par\smallskip{\small\color{muted}[KB09, BV18]}
\end{glossaryentry}


\begin{glossaryentry}{Matrix rank}
The dimension of the column space (equivalently, row space) of a matrix. Matrix rank should not be confused with CP rank or multilinear rank.

\end{glossaryentry}


\subsection{II. Decomposition models}

\begin{glossaryentry}{CP decomposition}
A model $\mathcal X\approx\sum_{r=1}^R\lambda_r a_r^{(1)}\otimes\cdots\otimes a_r^{(d)}$. CP is the main example where factor vectors are coordinates while the reconstructed tensor is the represented object.
\par\smallskip{\small\color{muted}[KB09, BV18]}
\end{glossaryentry}


\begin{glossaryentry}{CP rank}
The smallest number of rank-one tensors whose sum equals a tensor exactly. It differs from matrix rank and Tucker multilinear rank.
\par\smallskip{\small\color{muted}[KB09]}
\end{glossaryentry}


\begin{glossaryentry}{Factor matrix}
In CP, vectors from one mode are collected as columns of a factor matrix. Its columns are coordinates and can change under CP scaling or permutation equivalences.
\par\smallskip{\small\color{muted}[KB09]}
\end{glossaryentry}


\begin{glossaryentry}{Tucker decomposition}
A model $\mathcal X\approx\mathcal G\times_1U_1\times_2\cdots\times_dU_d$ with a core tensor and mode factors. Tucker emphasizes mode-wise subspaces and their interactions rather than a sum of rank-one terms.
\par\smallskip{\small\color{muted}[KB09, DLV00]}
\end{glossaryentry}


\begin{glossaryentry}{Core tensor}
The smaller tensor $\mathcal G$ in a Tucker representation that records interactions among selected mode coordinates. Individual core entries are basis-dependent; the mode subspaces can be more invariant than a particular basis.
\par\smallskip{\small\color{muted}[KB09]}
\end{glossaryentry}


\begin{glossaryentry}{Multilinear rank}
The tuple $(r_1,\ldots,r_d)$ of ranks of the tensor's mode unfoldings. It is the rank notion naturally associated with Tucker models.
\par\smallskip{\small\color{muted}[DLV00]}
\end{glossaryentry}


\begin{glossaryentry}{HOSVD}
Higher-Order SVD constructs Tucker mode subspaces from singular vectors of tensor unfoldings. It gives an operational way to inspect multilinear rank and compress a tensor.
\par\smallskip{\small\color{muted}[DLV00]}
\end{glossaryentry}


\begin{glossaryentry}{Block-Term Decomposition / BTD}
A sum of structured multilinear-rank blocks rather than only rank-one terms. A BTD component can represent within-component interactions that are not separable across modes.
\par\smallskip{\small\color{muted}[DL08]}
\end{glossaryentry}


\begin{glossaryentry}{SVD}
For a matrix $X$, $X=U\Sigma V^\top$. It provides an optimal low-rank approximation for a chosen matricization but can merge mode semantics after flattening.
\par\smallskip{\small\color{muted}[KB09]}
\end{glossaryentry}


\subsection{III. Geometry and equivalence}

\begin{glossaryentry}{Subspace}
A subset of a vector space closed under addition and scalar multiplication. Low-rank matrix and Tucker methods often identify a subspace more naturally than an individually meaningful basis vector.
\par\smallskip{\small\color{muted}[AMS08]}
\end{glossaryentry}


\begin{glossaryentry}{Orthonormal basis}
A basis whose vectors have unit norm and are mutually orthogonal. An orthonormal matrix is a convenient coordinate system for a subspace, but different bases can span the same subspace.

\end{glossaryentry}


\begin{glossaryentry}{Manifold}
A space that locally resembles Euclidean space and supports tangent-space calculus. Several low-rank model spaces are naturally treated as manifolds for optimization.
\par\smallskip{\small\color{muted}[AMS08]}
\end{glossaryentry}


\begin{glossaryentry}{Stiefel manifold}
$\mathrm{St}(n,r)=\{U\in\mathbb R^{n\times r}:U^\top U=I_r\}$, the set of matrices with orthonormal columns. It appears when learning orthonormal bases or subspace coordinates, as in StelLA.
\par\smallskip{\small\color{muted}[AMS08, LIS25]}
\end{glossaryentry}


\begin{glossaryentry}{Segre manifold / variety}
The geometric model of rank-one tensors obtained from outer products. A CP tensor is a sum of Segre components.
\par\smallskip{\small\color{muted}[BV18]}
\end{glossaryentry}


\begin{glossaryentry}{Fixed-rank manifold}
The smooth manifold of matrices having a fixed rank, away from rank-changing boundaries. It allows direct optimization of a low-rank matrix object without introducing a redundant factorization.
\par\smallskip{\small\color{muted}[AMS08, BIA25]}
\end{glossaryentry}


\begin{glossaryentry}{Coordinates vs represented object}
A parameterization gives coordinates to an object; different coordinates may represent exactly the same object. This is the central distinction of the tutorial: factors are coordinates, while the reconstructed matrix or tensor is the object.
\par\smallskip{\small\color{muted}[BV18]}
\end{glossaryentry}


\begin{glossaryentry}{Gauge freedom / equivalence}
A family of coordinate transformations that leaves the represented object unchanged.
\begin{interpretationnote}{Why it matters}
A change in factors does not necessarily mean a change in the represented object.
\end{interpretationnote}

\end{glossaryentry}


\begin{glossaryentry}{Scaling ambiguity}
A reciprocal rescaling of factors can preserve the same rank-one tensor, e.g. $(\alpha a)\otimes(\beta b)\otimes((\alpha\beta)^{-1}c)$. Factor norms can change substantially while reconstruction stays unchanged.
\par\smallskip{\small\color{muted}[KB09]}
\end{glossaryentry}


\begin{glossaryentry}{Permutation ambiguity}
Reordering CP components does not change their sum. Two CP outputs should be aligned before column-by-column comparison.
\par\smallskip{\small\color{muted}[KB09]}
\end{glossaryentry}


\begin{glossaryentry}{Identifiability}
A model is identifiable when the represented object determines its parameters uniquely after accepted equivalences are accounted for.
\begin{interpretationnote}{Interpretation boundary}
Identifiability supports coordinate recovery, but an identifiable factor is not automatically a scientifically meaningful concept.
\end{interpretationnote}
\par\smallskip{\small\color{muted}[KB09, DL08]}
\end{glossaryentry}


\subsection{IV. Optimization}

\begin{glossaryentry}{Nonconvex optimization}
Optimization in which the objective or feasible set does not have global convexity guarantees. Low-rank factorizations can have initialization-dependent trajectories and multiple stationary regions.

\end{glossaryentry}


\begin{glossaryentry}{Initialization}
The starting coordinates or starting low-rank object supplied to an iterative solver. Different starts can produce different optimization trajectories even for the same target and model.
\par\smallskip{\small\color{muted}[KB09]}
\end{glossaryentry}


\begin{glossaryentry}{ALS}
Alternating Least Squares updates one factor block at a time while keeping the others fixed. It is simple and effective, but its local least-squares systems can become poorly conditioned.
\par\smallskip{\small\color{muted}[KB09]}
\end{glossaryentry}


\begin{glossaryentry}{RALS}
Regularized ALS adds a stabilizing term to alternating updates. It illustrates that solver design can change behavior near difficult regions without changing the model family.

\end{glossaryentry}


\begin{glossaryentry}{Riemannian optimization}
Optimization performed on a manifold using tangent-space gradients and manifold-compatible updates. It respects the geometry of low-rank model spaces instead of treating all coordinates as unconstrained Euclidean parameters.
\par\smallskip{\small\color{muted}[AMS08]}
\end{glossaryentry}


\begin{glossaryentry}{Tangent space}
The linear space of infinitesimal feasible directions at a point on a manifold. Riemannian gradients and search directions live in tangent spaces.
\par\smallskip{\small\color{muted}[AMS08]}
\end{glossaryentry}


\begin{glossaryentry}{Retraction}
A map that sends a tangent-space step back to the manifold and agrees locally with the manifold geometry. It turns a tangent search direction into a new feasible point.
\par\smallskip{\small\color{muted}[AMS08]}
\end{glossaryentry}


\begin{glossaryentry}{RGD}
Riemannian Gradient Descent follows the negative Riemannian gradient and retracts each step. It is a basic geometry-aware alternative to Euclidean or alternating updates.
\par\smallskip{\small\color{muted}[AMS08]}
\end{glossaryentry}


\begin{glossaryentry}{RCG}
Riemannian Conjugate Gradient combines the current gradient with transported information from previous directions. It provides a geometry-aware acceleration strategy beyond steepest descent.
\par\smallskip{\small\color{muted}[AMS08]}
\end{glossaryentry}


\begin{glossaryentry}{Levenberg--Marquardt / LM}
A damped Gauss--Newton-type method for nonlinear least-squares problems. It provides a curvature-aware strategy for low-rank reconstruction.

\end{glossaryentry}


\subsection{V. Failure and stability}

\begin{glossaryentry}{Conditioning}
Sensitivity of a solution or representation to small perturbations of the input or represented object.
\begin{interpretationnote}{Interpretation boundary}
A low reconstruction error can coexist with component coordinates that are too sensitive for reliable interpretation.
\end{interpretationnote}
\par\smallskip{\small\color{muted}[BV18]}
\end{glossaryentry}


\begin{glossaryentry}{Condition number}
A numerical quantity that measures sensitivity; for nonsingular $M$, one common form is $\kappa(M)=\|M\|\|M^{-1}\|$. Different condition numbers diagnose different objects; ALS subproblem conditioning and geometric CP conditioning should not be conflated.
\par\smallskip{\small\color{muted}[BV18]}
\end{glossaryentry}


\begin{glossaryentry}{Relative reconstruction error}
$\|X-\widehat X\|_F/\|X\|_F$ (and analogously for tensors). It measures object-level fit, not uniqueness, stability, or semantic validity.
\par\smallskip{\small\color{muted}[KB09]}
\end{glossaryentry}


\begin{glossaryentry}{Overparameterization / excessive rank}
Using more latent components than are needed to represent the dominant structure. Extra components can split or duplicate structure and complicate interpretation.

\end{glossaryentry}


\begin{glossaryentry}{Rank choice}
The selected number of components or latent dimensions used by a model.
\begin{interpretationnote}{Interpretation boundary}
More rank can split or duplicate a vocabulary rather than merely improve the fit.
\end{interpretationnote}

\end{glossaryentry}


\begin{glossaryentry}{Component collision}
Two normalized rank-one components become nearly indistinguishable as represented rank-one tensors, up to accepted scale/sign conventions.
\begin{interpretationnote}{Why it matters}
Nearly indistinguishable rank-one terms can make component roles difficult to separate and local optimization problems poorly conditioned.
\end{interpretationnote}
\par\smallskip{\small\color{muted}[BV18]}
\end{glossaryentry}


\begin{glossaryentry}{Degeneracy}
A tensor-approximation pathology in which component norms can grow while their sum remains bounded, often through near cancellation. A small represented tensor can hide large, fragile component coordinates.
\par\smallskip{\small\color{muted}[KB09, BV18]}
\end{glossaryentry}


\begin{glossaryentry}{Cancellation}
Large terms with opposing contributions combine to produce a much smaller sum. It is a separate pathology from component collision, even though the two may occur together.

\end{glossaryentry}


\begin{glossaryentry}{Swamp / swamp-like stagnation}
An extended region of very slow objective improvement during an iterative tensor-decomposition method. A plateau is an observation, not an explanation: collision, degeneracy, initialization, rank choice, and model mismatch are possible contributors.
\par\smallskip{\small\color{muted}[KB09]}
\end{glossaryentry}


\begin{glossaryentry}{Model mismatch}
The chosen model family does not adequately represent the structure that generated the data. Poor fit is not always an optimization failure; the structural assumption itself may be wrong.

\end{glossaryentry}


\newpage
\section{AI and interpretability glossary}

\subsection{I. Neural representations}

\begin{glossaryentry}{Neural representation}
The internal activation state produced by a neural network for an input at a chosen layer or module. Low-rank methods are applied to collections of such states to test hypotheses about recurring structure.
\par\smallskip{\small\color{muted}[KIM18, GHO19, FEL23]}
\end{glossaryentry}


\begin{glossaryentry}{Activation}
The numerical output of a neuron, channel, token position, or other internal unit after a layer processes an input. In the CRAFT-style lab, each image crop produces an activation vector.
\par\smallskip{\small\color{muted}[FEL23]}
\end{glossaryentry}


\begin{glossaryentry}{Activation space}
The vector space in which a chosen layer's activation vectors live. CAVs, NMF directions, and SAE decoder directions can all be discussed as directions in an activation space.
\par\smallskip{\small\color{muted}[KIM18, CUN23]}
\end{glossaryentry}


\begin{glossaryentry}{Latent representation}
A lower-dimensional or internal representation that is not identical to the original input. CP, Tucker, NMF, SVD, and autoencoders all introduce latent coordinates under different assumptions.

\end{glossaryentry}


\begin{glossaryentry}{Representation learning}
Learning useful internal variables or coordinates from data rather than specifying all features by hand. The tutorial compares several structured representation families and asks what their coordinates can legitimately mean.

\end{glossaryentry}


\begin{glossaryentry}{Feature}
An overloaded term that may mean an input variable, neuron, activation direction, or learned latent factor. The tutorial avoids assuming that every learned factor is automatically a semantic feature.

\end{glossaryentry}


\begin{glossaryentry}{Feature direction}
A direction in activation space associated with a recurring learned pattern. Dictionary atoms, CAVs, and SAE decoder columns can all be treated as candidate feature directions, with different training assumptions.
\par\smallskip{\small\color{muted}[KIM18, CUN23]}
\end{glossaryentry}


\subsection{II. Dictionary representations}

\begin{glossaryentry}{Dictionary learning}
Learning a set of reusable basis elements $D$ such that observations can be approximated as combinations $x_i\approx D\alpha_i$. It is the broader representation-learning idea behind learned feature dictionaries; NMF and sparse autoencoders impose different additional structures.
\par\smallskip{\small\color{muted}[MBPS09]}
\end{glossaryentry}


\begin{glossaryentry}{Dictionary atom}
One learned basis element or direction in a dictionary.
\begin{interpretationnote}{Interpretation boundary}
A reusable learned direction is not automatically a human-meaningful concept.
\end{interpretationnote}
\par\smallskip{\small\color{muted}[MBPS09]}
\end{glossaryentry}


\begin{glossaryentry}{Sparse coding}
Representing each observation with only a small number of active dictionary coefficients. Sparsity and nonnegativity are distinct assumptions: a code can be sparse without being nonnegative and vice versa.
\par\smallskip{\small\color{muted}[MBPS09]}
\end{glossaryentry}


\begin{glossaryentry}{Sparsity}
The property that most entries of a representation are zero or inactive. Sparse representations encourage each observation to use only a small subset of learned features.
\par\smallskip{\small\color{muted}[MBPS09, CUN23]}
\end{glossaryentry}


\begin{glossaryentry}{Sparsity penalty}
A regularization term or constraint that encourages few active latent variables, such as an $\ell_1$ penalty or related activation penalty. In SAEs it trades reconstruction fidelity against sparse feature usage.

\end{glossaryentry}


\begin{glossaryentry}{Autoencoder}
A model trained to reconstruct an input through an intermediate latent code: $z=f_{enc}(x)$ and $\hat x=f_{dec}(z)$. It provides the basic architecture from which sparse autoencoders are built.

\end{glossaryentry}


\begin{glossaryentry}{Encoder}
The mapping $f_{enc}$ that converts an input or activation vector into latent coordinates. In an SAE, the encoder determines which learned latent features are active.
\par\smallskip{\small\color{muted}[CUN23]}
\end{glossaryentry}


\begin{glossaryentry}{Decoder}
The mapping $f_{dec}$ that reconstructs the input from the latent code. For a linear SAE decoder, decoder columns are commonly treated as learned feature directions or dictionary atoms.
\par\smallskip{\small\color{muted}[CUN23]}
\end{glossaryentry}


\begin{glossaryentry}{Latent code}
The intermediate representation $z$ produced by an encoder. In an SAE, sparsity is imposed on this code so that only a small number of features activate for each input.

\end{glossaryentry}


\begin{glossaryentry}{Sparse autoencoder / SAE}
An autoencoder trained to reconstruct activations through a latent code that is encouraged to be sparse.
\begin{interpretationnote}{Interpretation boundary}
An SAE feature is a learned representation direction, not automatically a validated semantic concept.
\end{interpretationnote}
\par\smallskip{\small\color{muted}[CUN23, BRI23, TEM26]}
\end{glossaryentry}


\begin{glossaryentry}{Overcomplete dictionary}
A dictionary with more learned atoms/features than the dimensionality of the input representation. Overcomplete feature sets are used to model the hypothesis that networks may represent more features than neurons/dimensions via superposition.
\par\smallskip{\small\color{muted}[CUN23, BRI23]}
\end{glossaryentry}


\begin{glossaryentry}{Feature activation}
The coefficient or latent value indicating how strongly a learned feature is active for one example. High activation means presence under that representation; it does not by itself establish importance or causal relevance.

\end{glossaryentry}


\begin{glossaryentry}{Dead feature / dead latent}
A learned latent feature that rarely or never activates on the data distribution used to evaluate it. Dead features indicate unused dictionary capacity and are a practical SAE training diagnostic.

\end{glossaryentry}


\begin{glossaryentry}{Monosemanticity}
The ideal that one learned feature corresponds to one coherent semantic pattern rather than several unrelated meanings. SAE work investigates whether sparse dictionaries yield more monosemantic features than individual neurons.
\par\smallskip{\small\color{muted}[CUN23, BRI23, TEM26]}
\end{glossaryentry}


\begin{glossaryentry}{Polysemanticity}
A neuron or learned feature responds to several semantically distinct patterns. It motivates attempts to recover cleaner latent directions than individual neurons provide.
\par\smallskip{\small\color{muted}[CUN23, BRI23]}
\end{glossaryentry}


\begin{glossaryentry}{Superposition}
The hypothesis that a network represents more features than available neurons/dimensions by encoding features in non-orthogonal directions. Sparse autoencoders and dictionary learning are used to search for those latent directions.
\par\smallskip{\small\color{muted}[CUN23, BRI23]}
\end{glossaryentry}


\begin{glossaryentry}{Feature splitting}
A phenomenon where a semantic pattern that might be described as one feature at one scale is represented by several more specific learned features at another scale or model capacity. It warns that increasing dictionary size can change the vocabulary rather than simply refine reconstruction.
\par\smallskip{\small\color{muted}[TEM26]}
\end{glossaryentry}


\begin{glossaryentry}{NMF}
A constrained factorization $A\approx UW^\top$ with $U,W\ge0$.
\begin{interpretationnote}{Interpretation boundary}
Nonnegativity encourages additive coordinates; it does not guarantee uniqueness or semantic meaning.
\end{interpretationnote}
\par\smallskip{\small\color{muted}[LS99, FEL23]}
\end{glossaryentry}


\subsection{III. Concept-based interpretability}

\begin{glossaryentry}{Concept}
A higher-level pattern used to describe model representations or behavior in human-relevant terms. A learned direction is first a candidate concept; semantic meaning requires additional evidence.
\par\smallskip{\small\color{muted}[KIM18, GHO19]}
\end{glossaryentry}


\begin{glossaryentry}{Concept-based interpretability}
Explaining neural networks using higher-level concepts rather than only individual input features or neurons. CRAFT is the main case study in Lab 4.
\par\smallskip{\small\color{muted}[KIM18, GHO19, FEL23]}
\end{glossaryentry}


\begin{glossaryentry}{Concept discovery}
Searching for recurring structure in neural activations that may correspond to meaningful higher-level patterns, often without pre-specifying all concepts. Discovery yields candidate explanations, not ground-truth semantics.
\par\smallskip{\small\color{muted}[GHO19, FEL23]}
\end{glossaryentry}


\begin{glossaryentry}{Feature attribution}
Assigning importance scores to input features or internal variables for a model prediction. Concept methods address a different level of explanation by grouping behavior into higher-level directions or concepts.
\par\smallskip{\small\color{muted}[GHO19, FEL23]}
\end{glossaryentry}


\begin{glossaryentry}{Concept attribution}
Estimating how a higher-level candidate concept contributes to a prediction. CRAFT combines concept discovery with concept importance and concept attribution maps.
\par\smallskip{\small\color{muted}[FEL23]}
\end{glossaryentry}


\begin{glossaryentry}{CAV}
A Concept Activation Vector: a direction in activation space associated with a concept.
\begin{interpretationnote}{Interpretation boundary}
A direction in activation space becomes a defensible concept only after semantic and behavioral validation.
\end{interpretationnote}
\par\smallskip{\small\color{muted}[KIM18, FEL23]}
\end{glossaryentry}


\begin{glossaryentry}{TCAV}
Testing with Concept Activation Vectors uses directional sensitivity in activation space to quantify relevance of user-defined concepts to predictions. It provides historical context for concept directions before automatic concept discovery methods such as ACE and CRAFT.
\par\smallskip{\small\color{muted}[KIM18]}
\end{glossaryentry}


\begin{glossaryentry}{ACE}
Automatic Concept-based Explanations automatically extracts visual concepts and evaluates their relevance to predictions. ACE is a key step from user-specified concepts toward automatic concept discovery.
\par\smallskip{\small\color{muted}[GHO19]}
\end{glossaryentry}


\begin{glossaryentry}{CRAFT}
Concept Recursive Activation FacTorization for Explainability: an automatic concept method using nonnegative activation factorization, recursive decomposition, Sobol-based importance, and concept attribution maps. It is the main interpretability case study in Lab 4.
\par\smallskip{\small\color{muted}[FEL23]}
\end{glossaryentry}


\begin{glossaryentry}{Image crop / patch}
A local region extracted from an image and analyzed as an example. CRAFT-style concept discovery uses crop activations so that candidate concepts can be associated with recurring visual parts or patterns.
\par\smallskip{\small\color{muted}[FEL23]}
\end{glossaryentry}


\begin{glossaryentry}{Global average pooling}
Averaging a convolutional activation map over spatial positions to obtain one feature value per channel. CRAFT uses pooled crop activations for its NMF concept-factorization stage.
\par\smallskip{\small\color{muted}[FEL23]}
\end{glossaryentry}


\begin{glossaryentry}{Concept bank}
A collection of candidate concept directions used to represent or analyze activations.
\begin{interpretationnote}{Interpretation boundary}
A collection of candidate directions is a learned representation, not yet a validated concept vocabulary.
\end{interpretationnote}
\par\smallskip{\small\color{muted}[FEL23]}
\end{glossaryentry}


\begin{glossaryentry}{Concept usage / coefficient}
A coefficient describing how strongly a candidate concept contributes to one observation's representation. In $A\approx UW^\top$, $U_{ij}$ is the usage of candidate $j$ by crop $i$.
\par\smallskip{\small\color{muted}[FEL23]}
\end{glossaryentry}


\begin{glossaryentry}{Concept importance}
How strongly model behavior depends on a candidate concept, rather than merely whether the concept is present.
\begin{interpretationnote}{Why it matters}
A concept can be strongly present in an activation while having little effect on the model's output.
\end{interpretationnote}
\par\smallskip{\small\color{muted}[FEL23]}
\end{glossaryentry}


\begin{glossaryentry}{Sobol sensitivity}
A variance-based global sensitivity measure that attributes output variation to input variables or groups. CRAFT applies Sobol indices to concept importance; simplified notebook meters should be labeled as pedagogical proxies unless they reproduce the full estimator.
\par\smallskip{\small\color{muted}[FEL23]}
\end{glossaryentry}


\begin{glossaryentry}{Concept Attribution Map}
A spatial map localizing evidence associated with a candidate concept in an input. It matters when comparing CRAFT's pooled NMF stage with tensor extensions that retain location directly.
\par\smallskip{\small\color{muted}[FEL23]}
\end{glossaryentry}


\begin{glossaryentry}{Recursive concept decomposition}
Re-analyzing a higher-level concept using activations from an earlier layer to obtain finer sub-concepts. It illustrates that the semantic granularity of a concept depends on the layer.
\par\smallskip{\small\color{muted}[FEL23]}
\end{glossaryentry}


\begin{glossaryentry}{Concept granularity}
The level of semantic detail expressed by a concept, from class-level structure to parts, textures, or other sub-concepts. Rank, layer, and representation family can all change the concept vocabulary and its granularity.
\par\smallskip{\small\color{muted}[FEL23]}
\end{glossaryentry}


\subsection{IV. Low-rank learning in neural networks}

\begin{glossaryentry}{LoRA}
Low-Rank Adaptation represents a trainable weight update with a low-rank parameterization rather than updating all pretrained weights. It motivates questions about whether learning factors, subspaces, or the fixed-rank object is the right geometric viewpoint.
\par\smallskip{\small\color{muted}[HUE22, LIS25]}
\end{glossaryentry}


\begin{glossaryentry}{Adapter}
A small trainable module or parameter update inserted into or associated with a pretrained model for parameter-efficient adaptation. LoRA is a low-rank adapter strategy.
\par\smallskip{\small\color{muted}[HUE22]}
\end{glossaryentry}


\begin{glossaryentry}{Low-rank weight update}
A weight change constrained to have low matrix rank. StelLA decomposes such an update as $USV^\top$ to make input/output subspaces explicit.
\par\smallskip{\small\color{muted}[LIS25]}
\end{glossaryentry}


\begin{glossaryentry}{Subspace learning}
Learning the low-dimensional space in which useful updates or representations live, rather than assigning meaning to one particular basis. StelLA explicitly learns orthonormal input and output subspaces using Stiefel factors.
\par\smallskip{\small\color{muted}[LIS25]}
\end{glossaryentry}


\begin{glossaryentry}{Factorized parameterization}
Representing a low-rank object through factors, e.g. $W=AB^\top$. It reduces parameter count but introduces non-unique coordinates and can affect conditioning.
\par\smallskip{\small\color{muted}[BIA25]}
\end{glossaryentry}


\begin{glossaryentry}{Parameter redundancy}
Different parameter values encode the same represented object or model behavior. RAdaGrad/RAdamW motivate direct fixed-rank optimization partly to avoid redundant factor coordinates.
\par\smallskip{\small\color{muted}[BIA25]}
\end{glossaryentry}


\begin{glossaryentry}{Fixed-rank weight}
A matrix parameter constrained to have a specified rank. It can be optimized as an intrinsic matrix object rather than only through a factorization.
\par\smallskip{\small\color{muted}[BIA25]}
\end{glossaryentry}


\begin{glossaryentry}{RAdaGrad / RAdamW}
Riemannian optimizers for fixed-rank matrix weights that adapt AdaGrad/AdamW-style metrics to the fixed-rank manifold. They illustrate geometry-aware low-rank learning without relying on a redundant factorization.
\par\smallskip{\small\color{muted}[BIA25]}
\end{glossaryentry}


\begin{glossaryentry}{Equivariance}
A property where transforming the input induces a predictable corresponding transformation of the representation or output. Tensor Decomposition Networks study low-rank approximations inside SO(3)-equivariant architectures.
\par\smallskip{\small\color{muted}[LIN25]}
\end{glossaryentry}


\begin{glossaryentry}{Tensor product}
A multilinear construction that combines representation spaces or features. Equivariant networks may use structured tensor-product interactions that can be expensive.
\par\smallskip{\small\color{muted}[LIN25]}
\end{glossaryentry}


\begin{glossaryentry}{Clebsch--Gordan tensor product}
A structured tensor-product operation used to combine SO(3)-equivariant features while respecting representation-theoretic symmetry. TDNs replace expensive CG tensor-product operators with low-rank tensor decompositions.
\par\smallskip{\small\color{muted}[LIN25]}
\end{glossaryentry}


\begin{glossaryentry}{Low-rank tensor operator}
A tensor-valued or multilinear operator approximated using a low-rank tensor decomposition. TDNs use CP structure to reduce the cost of tensor-product operations.
\par\smallskip{\small\color{muted}[LIN25]}
\end{glossaryentry}


\begin{glossaryentry}{Architectural compression}
Reducing computation or parameter cost by replacing a large operator or module with a structured lower-complexity approximation. In TDNs the low-rank tensor structure changes the architecture's operator implementation, not merely post-hoc interpretation.
\par\smallskip{\small\color{muted}[LIN25]}
\end{glossaryentry}


\begin{glossaryentry}{Separable interaction}
An interaction that factors into products of simpler mode-wise terms. CP decomposition imposes a separability assumption on each rank-one component.
\par\smallskip{\small\color{muted}[KB09, LIN25]}
\end{glossaryentry}


\subsection{V. Evidence and validation}

\begin{glossaryentry}{Mechanistic interpretability}
Studying internal model computations with the goal of identifying features, circuits, or mechanisms that explain behavior. SAE feature dictionaries are one current tool in this broader program.
\par\smallskip{\small\color{muted}[CUN23, BRI23]}
\end{glossaryentry}


\begin{glossaryentry}{Activation patching}
Replacing or restoring internal activations from another run to test whether a particular activation state contributes causally to behavior. It is an intervention-style diagnostic distinct from merely observing a large feature activation.

\end{glossaryentry}


\begin{glossaryentry}{Ablation}
Removing, zeroing, or disabling a feature, unit, component, or pathway and measuring the resulting behavioral change. It can provide causal evidence when appropriate controls rule out trivial distribution-shift effects.

\end{glossaryentry}


\begin{glossaryentry}{Intervention}
Deliberately changing a proposed mechanism or representation and testing whether the predicted effect follows. It is treated as stronger evidence than reconstruction or correlation alone.

\end{glossaryentry}


\begin{glossaryentry}{Causal relevance}
Evidence that changing a candidate feature or mechanism changes model behavior in the predicted way. A feature can be present or correlated without being causally important.
\par\smallskip{\small\color{muted}[CUN23]}
\end{glossaryentry}


\begin{glossaryentry}{Faithfulness}
The degree to which an explanation accurately reflects the model computation or behavior it claims to explain. CRAFT evaluates concept-importance faithfulness; visual plausibility alone is not enough.
\par\smallskip{\small\color{muted}[FEL23]}
\end{glossaryentry}


\begin{glossaryentry}{Semantic coherence}
The extent to which examples strongly associated with a learned feature support a consistent human interpretation. Coherence supports naming a candidate concept but does not alone establish causal importance.
\par\smallskip{\small\color{muted}[GHO19, FEL23]}
\end{glossaryentry}


\begin{glossaryentry}{Stability}
Persistence of a candidate structure across reasonable changes in initialization, samples, ranks, perturbations, or solver choices.
\begin{interpretationnote}{Interpretation boundary}
A candidate that does not persist across reasonable runs, ranks, samples, or perturbations is weak evidence for component-level interpretation.
\end{interpretationnote}

\end{glossaryentry}


\begin{glossaryentry}{Feature universality}
The hypothesis or observation that similar learned features recur across models, scales, seeds, or training settings. It can strengthen evidence that a feature is not a one-run artifact, but matching features across representations requires care.

\end{glossaryentry}


\begin{glossaryentry}{Held-out validation}
Testing a claim on data that were not used to fit or select the representation. A semantic label should generalize beyond the examples used to invent that label.

\end{glossaryentry}


\begin{glossaryentry}{Behavioral validation}
Testing whether a proposed concept is meaningfully related to model outputs on external examples or tasks. It separates 'this factor looks interpretable' from 'this factor predicts something about model behavior'.
\par\smallskip{\small\color{muted}[KIM18, GHO19, FEL23]}
\end{glossaryentry}


\begin{glossaryentry}{Semantic validation}
Evidence that a mathematical factor corresponds to the domain meaning assigned to it. High-usage examples can suggest a label; stability, held-out behavior, or interventions strengthen the claim.

\end{glossaryentry}


\begin{glossaryentry}{Reconstruction}
How accurately a learned representation reproduces the data or activation object it was fit to.
\begin{interpretationnote}{Interpretation boundary}
Good reconstruction does not establish uniqueness, stability, semantic meaning, or causal relevance.
\end{interpretationnote}

\end{glossaryentry}


\begin{glossaryentry}{Alignment-relevant evidence}
Representation- or behavior-level evidence that may support auditing or monitoring questions relevant to system behavior. The tutorial does not present tensor decomposition as an alignment algorithm; it studies when low-rank evidence is sufficiently invariant, stable, and behaviorally validated to be useful for auditing.

\end{glossaryentry}


\newpage
\section{Translation guide: mathematics to interpretability}
\begin{longtable}{p{0.23\linewidth} p{0.28\linewidth} p{0.41\linewidth}}
\toprule
\textbf{Mathematical language} & \textbf{Interpretability language} & \textbf{Caution}\\
\midrule
Factor / dictionary atom & Candidate feature or concept direction & A mathematical direction becomes semantic only after validation. \\
Coefficient & Feature or concept usage & Presence is not importance. \\
Subspace & Representation/update subspace & A stable subspace need not have a unique meaningful basis. \\
Rank & Vocabulary capacity / granularity & More rank can split or duplicate concepts. \\
Conditioning & Representation stability & Good reconstruction can coexist with fragile coordinates. \\
Identifiability & Recoverability & Recoverability is not domain meaning. \\
Perturbation / intervention & Behavioral evidence & The claim depends on the intervention and its controls. \\
\bottomrule
\end{longtable}

\keydistinction{A mathematically recoverable factor is not automatically a scientifically meaningful concept.}

\section{Reference index}
\begin{description}[leftmargin=1.55cm,style=nextline,itemsep=.55em]
\item[\textbf{[KB09]}] T. G. Kolda and B. W. Bader. Tensor Decompositions and Applications. SIAM Review 51(3):455--500, 2009. DOI 10.1137/07070111X.
\item[\textbf{[DLV00]}] L. De Lathauwer, B. De Moor, and J. Vandewalle. A Multilinear Singular Value Decomposition. SIAM Journal on Matrix Analysis and Applications 21(4):1253--1278, 2000.
\item[\textbf{[DL08]}] L. De Lathauwer. Decompositions of a Higher-Order Tensor in Block Terms---Part II: Definitions and Uniqueness. SIAM Journal on Matrix Analysis and Applications 30(3):1033--1066, 2008. DOI 10.1137/070690729.
\item[\textbf{[AMS08]}] P.-A. Absil, R. Mahony, and R. Sepulchre. Optimization Algorithms on Matrix Manifolds. Princeton University Press, 2008.
\item[\textbf{[BV18]}] P. Breiding and N. Vannieuwenhoven. A Riemannian Trust Region Method for the Canonical Tensor Rank Approximation Problem. SIAM Journal on Optimization 28(3):2435--2465, 2018. DOI 10.1137/17M114618X.
\item[\textbf{[LS99]}] D. D. Lee and H. S. Seung. Learning the Parts of Objects by Non-negative Matrix Factorization. Nature 401:788--791, 1999. DOI 10.1038/44565.
\item[\textbf{[MBPS09]}] J. Mairal, F. Bach, J. Ponce, and G. Sapiro. Online Dictionary Learning for Sparse Coding. ICML 2009, pp. 689--696. DOI 10.1145/1553374.1553463.
\item[\textbf{[GBC16]}] I. Goodfellow, Y. Bengio, and A. Courville. Deep Learning. MIT Press, 2016.
\item[\textbf{[KIM18]}] B. Kim et al. Interpretability Beyond Feature Attribution: Quantitative Testing with Concept Activation Vectors (TCAV). ICML 2018, PMLR 80:2668--2677.
\item[\textbf{[GHO19]}] A. Ghorbani, J. Wexler, J. Y. Zou, and B. Kim. Towards Automatic Concept-based Explanations. NeurIPS 2019, pp. 9277--9286.
\item[\textbf{[FEL23]}] T. Fel et al. CRAFT: Concept Recursive Activation FacTorization for Explainability. CVPR 2023, pp. 2711--2721.
\item[\textbf{[HUE22]}] E. J. Hu et al. LoRA: Low-Rank Adaptation of Large Language Models. ICLR 2022.
\item[\textbf{[CUN23]}] H. Cunningham, A. Ewart, L. Riggs, R. Huben, and L. Sharkey. Sparse Autoencoders Find Highly Interpretable Features in Language Models. arXiv:2309.08600, 2023.
\item[\textbf{[BRI23]}] T. Bricken et al. Towards Monosemanticity: Decomposing Language Models With Dictionary Learning. Anthropic research report, 2023.
\item[\textbf{[TEM26]}] A. Templeton et al. Scaling Monosemanticity: Extracting Interpretable Features from Claude 3 Sonnet. arXiv:2605.29358, 2026.
\item[\textbf{[LIS25]}] Z. Li, S. Sajadmanesh, J. Li, and L. Lyu. StelLA: Subspace Learning in Low-rank Adaptation using Stiefel Manifold. NeurIPS 38, 2025.
\item[\textbf{[BIA25]}] F. Bian, J. Zheng, Z. Liu, J. Luo, and J.-F. Cai. Finding Low-Rank Matrix Weights in DNNs via Riemannian Optimization: RAdaGrad and RAdamW. NeurIPS 38, 2025.
\item[\textbf{[LIN25]}] Y. Lin et al. Tensor Decomposition Networks for Fast Machine Learning Interatomic Potential Computations. NeurIPS 38, 2025.
\end{description}


\end{document}
